\section{Introduction}
\label{sec:introduction}

The Voting Rights Act's Section 2 obligation to create majority-minority districts---where minority communities are sufficiently large, geographically compact, and politically cohesive to support such a district---has been litigated almost exclusively in the context of congressional redistricting. The seminal cases (\textit{Thornburg v.\ Gingles}, \textit{Shaw v.\ Reno}, \textit{Miller v.\ Johnson}, \textit{Allen v.\ Milligan}, and \textit{Louisiana v.\ Callais}) all address congressional maps drawn to the scale of one district per 764,000 persons.\footnote{The average congressional district population in 2020 is approximately 764,000 persons (total U.S. population 331.4 million divided by 435 House seats).}

State legislative redistricting operates at a different scale. State house districts range from 14,000 persons in New Hampshire (with 400 districts) to 146,000 persons in California (with 80 districts), with a national mean of approximately 47,000 persons per state house district. This scale difference has three implications for VRA analysis. First, more minority opportunity districts can potentially be created at smaller scale, because a minority population concentrated in a geographic area of a given size may be sufficient to constitute a majority in a small state house district but not in a large congressional district. Second, the geographic compactness requirement for majority-minority districts (the first \textit{Gingles} precondition) may be easier to satisfy when districts are smaller, because a smaller district can be drawn to more precisely match the geographic extent of a concentrated minority community. Third, the Callais disentanglement problem---separating race-motivated from partisan-motivated districting---applies at state legislative scale with the same constitutional constraints as at congressional scale.

This paper addresses three questions about VRA compliance in algorithmic state legislative redistricting. First, does the 42\% minority population threshold identified as a practical bright line for majority-minority district formation at the congressional level \citep{dellaLibera2026threshold} hold at state house scale? Second, does the Callais disentanglement methodology---which establishes that race-neutral algorithmic maps can satisfy Section 2 without running afoul of the Equal Protection Clause's anti-racial-classification principle---apply at state legislative scale? Third, how do the numbers of minority opportunity districts achievable through algorithmic state house redistricting compare to those achievable through algorithmic congressional redistricting in the same states?

Our findings are encouraging. The 42\% threshold holds at state legislative scale: VRASection applied to state house chambers in the five covered states produces majority-minority districts at the same threshold calibration as congressional-level analysis. The Callais disentanglement applies identically at state scale. And state house maps achieve minority opportunity district targets in all 5 covered states at state house scale---exceeding the congressional success rate (4 of 5)---with additional precision because the smaller district scale allows more exact matching of minority community boundaries. The 42\% threshold applies at state house scale in the same manner as at congressional scale (D.1).

The paper proceeds as follows. Section~\ref{sec:methodology} describes the VRASection algorithm and its application to state house chambers. Section~\ref{sec:threshold} analyzes the 42\% threshold at state house scale. Section~\ref{sec:callais} examines the Callais disentanglement at state legislative scale. Section~\ref{sec:comparison} compares minority opportunity district creation across scales. Section~\ref{sec:conclusion} concludes.

\subsection{Legal Background: Section 2 at State Scale}

Section 2 of the Voting Rights Act prohibits voting practices or procedures that ``result in a denial or abridgement of the right of any citizen of the United States to vote on account of race or color.''\footnote{52 U.S.C. \S\,10301(a).} Under \textit{Thornburg v.\ Gingles}, a Section 2 vote dilution claim requires showing that: (1) the minority community is sufficiently large and geographically compact to constitute a majority in a single-member district; (2) the minority community is politically cohesive; and (3) the majority typically votes as a bloc in a manner that defeats the minority's preferred candidate.\footnote{\textit{Thornburg v.\ Gingles}, 478 U.S. 30, 50--51 (1986).}

These preconditions apply to \textit{any} single-member district system, including state legislative districts. \textit{Reynolds v.\ Sims} established the one-person-one-vote requirement for state legislative districts;\footnote{\textit{Reynolds v.\ Sims}, 377 U.S. 533 (1964).} the VRA's Section 2 protections apply to state legislative elections on equal terms with congressional elections. There is no doctrinal basis for applying different standards to congressional and state legislative redistricting under Section 2.

In practice, however, state legislative Section 2 litigation has been less common than congressional litigation, in part because the smaller district sizes make the geographic compactness precondition easier to satisfy. Our analysis suggests that this practical pattern may also reflect the fact that Section 2 obligations at state legislative scale are, if anything, \textit{easier} to satisfy than at congressional scale---more districts mean more opportunities to create majority-minority configurations without excessive departure from compactness standards.
